\section{Can Love Be Commanded?}

Most of us know the difference between doing the right thing and wanting the good we are doing. A promise can be kept resentfully. The truth can be told while we wish it made no claim upon us. We can kneel to pray while the heart hangs back. Duty can govern the act before it has won the will.

Scripture does something startling: it commands love. Israel is to love the Lord with the whole heart, soul, and strength; Christ names this the first and greatest commandment and joins it to love of neighbor (Deut 6:4--5; Matt 22:37--40). Love is not left as an elective refinement for those who have completed the serious work of obedience. The law itself is ordered toward it.

That demand is unsettling. Conduct can be compelled: words repeated, tribute paid, attendance enforced, injury restrained. But love produced only by force would not be love. If love must be free, how can it be owed? If it is owed, how can it remain free?

Freedom is not the absence of a true end, as though the will were freest when nothing deserved its love. And a command is not a machine for manufacturing the inward act it names. A person may be bound to tell the truth while lying willingly; obligation identifies the good without causing the will to embrace it. The command to love is more searching still. It addresses the principle from which the rest of moral life proceeds.

Aquinas confronts the difficulty directly. A precept concerns what is due, and union with God is the end of spiritual life. Charity, by which that union begins, may therefore be commanded. The obligation of the command is opposed to liberty only in a will averse from what is commanded. Love itself can be fulfilled only willingly (ST II--II, q.~44, a.~1, especially ad 2).

Before going further, the word \emph{love} must be kept exact. Charity is not intense feeling, indulgence, attraction, or an affectionate name for whatever the will already prefers. It is the theological virtue by which God is loved above all things for God's own sake and the neighbor is loved in relation to God (ST II--II, q.~23, aa.~1 and 5; CCC 1822). Feeling may accompany charity, but neither emotional warmth nor its absence settles whether charity is present. The command reaches the will's fundamental direction.

Love may be commanded. It cannot be manufactured by an external command.
